% Tables fixture for ReadableSTEM (written for the tests; no third-party text).
% Data tables in the styles LaTeX users write (booktabs, a ruled grid, no rules,
% multirow/multicolumn headers, maths in cells, a blank worksheet table) and things
% that must not become tables (a table of contents, a two-column grid of prose,
% matrices, an augmented array, align and cases). tests/test_tables_fixtures.py.
\documentclass[11pt]{article}
\usepackage{amsmath}
\usepackage{booktabs}
\usepackage{multirow}
\usepackage{array}
\usepackage[margin=1in]{geometry}
\pagestyle{plain}
\begin{document}

\tableofcontents

\section{Free fall}
\subsection{Drop times}

Each group dropped a steel ball from five heights and timed the fall with a photogate.
The acceleration follows from the height and the time of the fall.

\begin{table}[h]
\centering
\caption{Drop times for five heights.}
\begin{tabular}{ccc}
\toprule
Height (m) & Time (s) & $g$ (m/s$^2$) \\
\midrule
0.50 & 0.319 & 9.83 \\
0.75 & 0.391 & 9.81 \\
1.00 & 0.452 & 9.79 \\
1.25 & 0.505 & 9.80 \\
1.50 & 0.553 & 9.81 \\
\bottomrule
\end{tabular}
\end{table}

The values agree within half a per cent, which is less than the uncertainty of a
single timing.

\subsection{Quiz scores}

\begin{table}[h]
\centering
\caption{Frequency of quiz scores.}
\begin{tabular}{|c|c|c|}
\hline
Score & Students & Share \\
\hline
6 & 3 & 0.15 \\
\hline
7 & 6 & 0.30 \\
\hline
8 & 7 & 0.35 \\
\hline
9 & 4 & 0.20 \\
\hline
\end{tabular}
\end{table}

\clearpage
\section{Tables without rules}
\subsection{Values of a cubic}

The table below gives $f(x) = x^3 - 2x$ at evenly spaced points. It is set without any
rules, as many instructors type tables in a hurry.

\begin{table}[h]
\centering
\begin{tabular}{rrrr}
$x$ & $f(x)$ & $f'(x)$ & Difference \\[4pt]
$-2$ & $-4$ & $10$ & $5.0$ \\
$-1$ & $1$ & $1$ & $-1.0$ \\
$0$ & $0$ & $-2$ & $-1.0$ \\
$1$ & $-1$ & $1$ & $5.0$ \\
$2$ & $4$ & $10$ & $17.0$ \\
\end{tabular}
\caption{Values of $f$ and its derivative.}
\end{table}

\subsection{Heat of materials}

\begin{table}[h]
\centering
\caption{Specific heat by phase.}
\begin{tabular}{llcc}
\toprule
\multirow{2}{*}{Type} & \multirow{2}{*}{Material} & \multicolumn{2}{c}{Specific heat (J/kg K)} \\
\cmidrule(lr){3-4}
 & & Solid & Liquid \\
\midrule
\multirow{2}{*}{Metal} & Aluminium & 900 & 1180 \\
 & Copper & 385 & 530 \\
\midrule
\multirow{2}{*}{Other} & Water & 2090 & 4186 \\
 & Ethanol & 2420 & 2440 \\
\bottomrule
\end{tabular}
\end{table}

\clearpage
\section{Quantities and a worksheet}
\subsection{Definitions}

\begin{table}[h]
\centering
\caption{Mechanical quantities.}
\renewcommand{\arraystretch}{1.6}
\begin{tabular}{llc}
\toprule
Quantity & Symbol & Definition \\
\midrule
Momentum & $\vec{p}$ & $\vec{p} = m\vec{v}$ \\
Kinetic energy & $K$ & $K = \tfrac{1}{2}mv^2$ \\
Work & $W$ & $W = \vec{F}\cdot\vec{d}$ \\
Power & $P$ & $P = W/t$ \\
\bottomrule
\end{tabular}
\end{table}

\subsection{Record your data}

Fill in one row for each trial during the lab.

\begin{table}[h]
\centering
\caption{Pendulum trials.}
\renewcommand{\arraystretch}{1.8}
\begin{tabular}{|p{2cm}|p{3cm}|p{3cm}|p{3cm}|}
\hline
Trial & Length (m) & Ten swings (s) & Period (s) \\
\hline
 & & & \\
\hline
 & & & \\
\hline
 & & & \\
\hline
 & & & \\
\hline
\end{tabular}
\end{table}

\clearpage
\section{Not tables}
\subsection{Key terms}

\noindent
\begin{tabular}{@{}p{3cm}p{11cm}@{}}
\textbf{Displacement} & The change in position of an object, measured as a straight line from where it started to where it ended, with a direction. \\[6pt]
\textbf{Velocity} & How fast the position changes and in which direction; the slope of a position against time graph at any instant. \\[6pt]
\textbf{Acceleration} & How fast the velocity changes; it can point against the motion, in which case the object slows down. \\
\end{tabular}

\subsection{Matrices and systems}

The rotation by a quarter turn and the system it solves are
\[
R = \begin{pmatrix} 0 & -1 \\ 1 & 0 \end{pmatrix}, \qquad
A = \begin{pmatrix} 2 & 1 & 0 \\ 1 & 3 & 1 \\ 0 & 1 & 4 \end{pmatrix}, \qquad
\left[\begin{array}{ccc|c} 2 & 1 & 0 & 3 \\ 1 & 3 & 1 & 5 \\ 0 & 1 & 4 & 5 \end{array}\right].
\]
Elimination gives
\begin{align*}
2x + y &= 3, \\
x + 3y + z &= 5, \\
y + 4z &= 5,
\end{align*}
and the sign function is
\[
\operatorname{sgn}(x) = \begin{cases} -1 & \text{if } x < 0, \\ 0 & \text{if } x = 0, \\ 1 & \text{if } x > 0. \end{cases}
\]

\end{document}
